\bloque{metodos-3d}
\tema{coordenadas}{Coordenadas en el espacio}
\dificultad{1}
\section{Puntos y ejes del espacio}
Un punto del espacio se describe mediante tres coordenadas respecto a unos ejes
de referencia. En un sistema cartesiano los ejes son perpendiculares entre sí.
Cada coordenada indica un desplazamiento con signo desde el origen en la dirección
del eje correspondiente. El orden de las coordenadas es relevante y debe mantenerse
durante todos los cálculos y al interpretar una representación gráfica.
El punto $(1,2,3)$ se sitúa usando una unidad en el primer eje,
dos unidades en el segundo y tres unidades en el tercero. La posición aparente
en la pantalla depende de la cámara y no modifica las coordenadas matemáticas.
Al girar la cámara observamos el mismo objeto desde otra dirección.
Al acercarla aumenta su tamaño aparente, sin cambiar las distancias del modelo.
Para orientarse conviene localizar el origen y seguir los colores de los ejes.
Una proyección plana puede hacer coincidir visualmente puntos diferentes del espacio.
Por eso es útil comparar varias perspectivas antes de decidir qué relación existe
entre dos objetos. Describe las coordenadas y separa los datos del problema de
las impresiones visuales que produce una perspectiva concreta.
